\documentclass[10pt,a4]{examenEPN0}
\usepackage{comandosEPN}

\usepackage{empheq}

\newcommand{\itemcp}{\stepcounter{enumi}\item[ \textcolor{red}{\theenumi.$^*$}] }

\newcommand{\iitemcp}{\stepcounter{enumii}\item[ \textcolor{red}{\theenumii)$^*$}] }


\materia{Álgebra Lineal}
%\autor{Departamento de Formación Básica}
\autor{GR \hspace{1cm}.}
%\fecha{Semestre 2026-A}
\fecha{Nombre: \hspace{9cm} Cédula:}
\logouno[1.6cm]{Logo_EPN.eps}
\logodos[2.6cm]{Logo_DFB.eps}

\newcommand{\la}{\lambda}

%Datos
\newcommand{\datosAlgebra}[2]{
\EnExamen
{
\vspace{5mm}
    \noindent
    Nombres: \hrulefill\
    Carrera: \rule{4.5cm}{0.4pt}\
    \ifthenelse{\equal{#1}{}}%
    {Paralelo: \underline{\parbox{1.8cm}{\parbox{1cm}{}}}\\[3mm]}
    {Paralelo: \underline{\parbox{1.8cm}{\centering #1}}\\[3mm]}
    C\'edula: \rule{5cm}{0.4pt}\
    Firma: \hrulefill\
    Fecha: \underline{\parbox{2.1cm}{\centering #2}}
}}


% \newcommand{\indicacionesPrueba}{%
% \EnExamen{
% \begin{indicaciones}
%     \item La prueba tiene una duraci\'on  de \textbf{120 minutos.}
%     \item No se permite la utilizaci\'on de dispositivos electr\'onicos (celular, tablets, calculadora, etc.)
%     \item Leer detenidamente los enunciados, \textbf{justificar la respuesta para acceder al puntaje total}.
%     \item La prueba es completamente individual. Los estudiantes deberán dejar sus teléfonos celulares guardados y no podrán tenerlos consigo (por ejemplo, en los bolsillos de su ropa) durante toda la evaluación. En caso de detectarse algún acto de deshonestidad académica o el incumplimiento de esta norma, no podrán seguir rindiendo la prueba y su calificación será \textbf{cero}.
% \end{indicaciones}
% \pagestyle{empty}
% }}

\newcommand{\indicacionesPrueba}{%
\EnExamen{
\begin{indicaciones}
    \item La prueba tiene una duraci\'on  de \textbf{60 minutos.}
    \item Leer detenidamente los enunciados, \textbf{justificar la respuesta} para acceder al puntaje total.
    \item No se permite la utilizaci\'on de dispositivos electr\'onicos (celular, tablets, calculadora, etc.)
    \item La prueba es completamente individual
    \item No se permite portar teléfonos celulares, relojes inteligentes ni otros dispositivos electrónicos durante el examen, aun cuando permanezcan en los bolsillos. Estos dispositivos deberán guardarse en la mochila antes del inicio de la evaluación y esta deberá mantenerse cerrada durante todo el examen.
    \item He leído las indicaciones y acepto que el incumplimiento de estas se penaliza con \textbf{cero}. Firma:...................................

    
\end{indicaciones}
\pagestyle{empty}
}}

\newcommand{\indicacionesExamen}{
\EnExamen{\begin{indicaciones}
    \item La evaluación tiene una duraci\'on  de \textbf{120 minutos.}
    \item No se permite la utilizaci\'on de dispositivos electr\'onicos (celular, tablets, calculadora, etc.)
    %\item Leer detenidamente los enunciados, justificar la respuesta para acceder al puntaje total.
    \item El examen es completamente individual. En caso de detectarse alg\'un acto de deshonestidad acad\'emica no podr\'a seguir rindiendo la prueba y su calificaci\'on ser\'a \textbf{cero}.
\end{indicaciones}
}\pagestyle{empty}
}

%% Imprimir las una hoja A4, dos veces en A5
\newcommand{\imprimirDoble}
{%
\usepackage{pgfmorepages}
\pgfmorepagesloadextralayouts
\pgfpagesuselayout{repeated 2-up}[a4paper,landscape]
\thispagestyle{empty}
}


%Este sí se puede comentar
\newenvironment{solucion}{\color{blue}\proof[Solución]}{\endproof\color{black}}




% \renewcommand{\@autor}{
% ---\\
% \phantom{.} 
% ...}
% \makeatother

\examen{Examen 01: solucionario y rúbrica}
\makeatletter
\renewcommand{\@autor}{Departamento de Formación Básica}
\renewcommand{\@fecha}{Miércoles de 7 a 9 -- 18/11/2026}
\makeatother
\begin{document}
\encabezado
\noindent\textbf{Criterio general (20 puntos).} Se asigna crédito parcial por
los pasos correctos. Un error aritmético se penaliza donde se origina;
los pasos posteriores se valoran por su coherencia, sin repetir la penalización.
Se aceptan procedimientos equivalentes debidamente justificados.

\section*{1. Matriz con parámetro (4 puntos)}
\noindent\textbf{Rúbrica.}
\begin{itemize}[itemsep=2pt]
    \item Calcular el determinante con procedimiento: 1,5 puntos; establecer la condición de invertibilidad: 0,5 puntos.
    \item Sustituir $\alpha=0$ y reducir el sistema: 1 punto; expresar todas las soluciones con una variable libre: 1 punto.
\end{itemize}
Aplicando $F_2\leftarrow F_2-2F_1$ y
$F_3\leftarrow F_3-3F_1$, el determinante es
\[
    \det(A)=\det\begin{pmatrix}1&\alpha&2\\0&\alpha&2\\0&\alpha&\alpha-4\end{pmatrix}
    =\alpha(\alpha-4)-2\alpha=\alpha(\alpha-6).
\]
Por tanto, $A$ es invertible si y solo si $\alpha\notin\{0,6\}$.
Para $\alpha=0$, las ecuaciones son $x+2z=1$, $2x+6z=2$ y
$3x+2z=3$. Restando dos veces la primera de la segunda se obtiene
$2z=0$; así $x=1$ y $y$ es libre. La tercera ecuación se satisface.
El conjunto solución es $\boxed{\{(1,t,0):t\in\R\}}$.

\section*{2. Intersección de subespacios de matrices (4 puntos)}
\noindent\textbf{Rúbrica.}
\begin{itemize}[itemsep=2pt]
    \item Imponer simultáneamente las condiciones de $S$ y $T$: 1 punto; parametrizar la intersección: 1 punto.
    \item Plantear las condiciones sobre $k$: 1 punto; hallar el valor y verificar todas las condiciones: 1 punto.
\end{itemize}
En $T$ las matrices tienen la forma $\begin{pmatrix}a&b\\b&a\end{pmatrix}$.
La condición de $S$ exige $2a+2b=0$, luego
\[
    \boxed{S\cap T=\left\{t\begin{pmatrix}1&-1\\-1&1\end{pmatrix}:t\in\R\right\}.}
\]
Para la matriz dada, $k+1=3k-1$ y $k-3=2k-4$ dan $k=1$.
La suma de entradas es $7k-7$, que también se anula para ese valor.
La matriz resultante es $\begin{pmatrix}2&-2\\-2&2\end{pmatrix}$.

\section*{3. Espacio generado (4 puntos)}
\noindent\textbf{Rúbrica.}
\begin{itemize}[itemsep=2pt]
    \item Formar la matriz de generadores: 0,5 puntos; calcular su determinante: 1 punto; concluir los valores de $k$: 0,5 puntos.
    \item Para $k=2$, describir el plano: 0,5 puntos; justificar su dimensión y que no es $\R^3$: 0,5 puntos.
    \item Dar un vector permitido del plano y justificarlo: 0,5 puntos; dar un vector exterior y justificarlo: 0,5 puntos.
\end{itemize}
Se tiene $u\times v=(5,-3,-1)$ y
$\det[u\ v\ w]=(u\times v)\cdot w=2-k$.
Los tres vectores generan $\R^3$ exactamente cuando $k\ne2$.
Si $k=2$, $w=v-u$, y $u,v$ son independientes (no son proporcionales).
Por tanto,
\[
    \operatorname{span}\{u,v,w\}=\{(x,y,z)\in\R^3:5x-3y-z=0\}.
\]
Es un plano por el origen, de dimensión 2, y no todo $\R^3$.
Por ejemplo, $u+v=(3,5,0)$ pertenece al plano y es distinto de cero y de
los tres generadores. El vector $(1,0,0)$ no pertenece porque $5\ne0$.

\section*{4. Vehículos (4 puntos)}
\noindent\textbf{Rúbrica.}
\begin{itemize}[itemsep=2pt]
    \item Definir las variables: 0,25 puntos; plantear las tres ecuaciones: 0,75 puntos (0,25 por ecuación); identificar $A$, $\mathbf{x}$ y $b$: 1 punto.
    \item Resolver el sistema con procedimiento: 1,5 puntos; interpretar y comprobar las cantidades: 0,5 puntos.
\end{itemize}
Sean $a,b,c$ las cantidades de vehículos A, B y C, enteras no negativas.
El modelo y su representación matricial son
\[
    a+b+c=9,\qquad a+2b+3c=17,\qquad 2a+b+c=12,
\]
\[
    A=\begin{pmatrix}1&1&1\\1&2&3\\2&1&1\end{pmatrix},\quad
    \mathbf{x}=\begin{pmatrix}a\\b\\c\end{pmatrix},\quad
    b_{\rm modelo}=\begin{pmatrix}9\\17\\12\end{pmatrix}.
\]
Restando la primera ecuación de la tercera, $a=3$. Entonces $b+c=6$
y $2b+3c=14$, de donde $c=2$ y $b=4$.
Salieron \textbf{3 vehículos A, 4 B y 2 C}:
$3+4+2=9$, $3+8+6=17$ y $6+4+2=12$.

\section*{5. Subespacio de polinomios (4 puntos)}
\noindent\textbf{Rúbrica.}
\begin{itemize}[itemsep=2pt]
    \item Verificar que el polinomio cero pertenece a $W$: 1 punto.
    \item Demostrar la cerradura bajo suma usando ambas condiciones: 1,5 puntos.
    \item Demostrar la cerradura bajo producto por escalar usando ambas condiciones: 1 punto.
    \item Concluir que es subespacio: 0,5 puntos.
\end{itemize}
El polinomio cero satisface $p(2)=p'(0)=0$. Si $p,q\in W$, entonces
$(p+q)(2)=p(2)+q(2)=0$ y $(p+q)'(0)=p'(0)+q'(0)=0$.
Para $\lambda\in\R$, $(\lambda p)(2)=\lambda p(2)=0$ y
$(\lambda p)'(0)=\lambda p'(0)=0$. La suma y el producto por escalar
permanecen en $\mathcal P_3(\R)$, luego $W$ es un subespacio.
También se acepta justificarlo como intersección de los núcleos de las
aplicaciones lineales $p\mapsto p(2)$ y $p\mapsto p'(0)$.
\end{document}
